\documentclass{article}
\usepackage{amsmath}
\usepackage{amssymb}
\usepackage[margin=1in]{geometry}
\newcommand{\sectionbreak}{\clearpage}

\title{LaTeX Test Cases}
\author{Joe Ricks}
\date{\today}

\begin{document}

\maketitle

\section{Problem 1}

Translate the tablets below and answer the questions on the next page.
Submit just the second page with your solutions

\begin{itemize}
    \item Describe your process of solving in detail.
    
    In order to solve this problem, we first needed to make some assumptions (or propositions as referred to in Rosen Section 1.1) about the tablet.

    1. We assume that the symbols represent numbers in a base-10 system. 

    2. We assume the numbers displayed side by side is representative of a relationship between them (i.e. there is a relationship between each pair).

    3. The symbol on the first line represents the number 1.
    With these assumptions, we proceeded to attempt to decode the symbols into standard base-10 numbers. We represet each of the pairs in tuples, defined below:
    \begin{align*}
        \\T = \{(1, 9), (2, 18), (3, 27), (4, 36),
        \\(5, 45), (6, 54), (7, 63), (8, 72), (9, 81), \\(10, 90), (11, 99), (12, 108), (13, 117), (14, 126)\}
    \end{align*}
    From here, we noticed that the second number in each tuple is simply 9 times the first number. This led us to conclude that the relationship represented by the symbols is a multiplication by 9.
    \item Are you using inductive or deductive reasoning?
    
    Per the definitions provided in the workbook, inductive reasoning is the method of drawing general conclusions from a limited set of observations, while deductive reasoning is the method of using logic to draw conclusions from statements we
accept as true. We are using inductive reasoning. We observed specific instances (the pairs of numbers) and identified a pattern (multiplication by 9) that we generalized to form a conclusion about the relationship represented by the symbols.
    
    \item Can you be sure your answer is correct?
    
    While we cannot be absolutely certain without additional context or confirmation, the consistent pattern observed across all pairs strongly suggests that our conclusion is correct. The simplicity and uniformity of the relationship (multiplication by 9) across all provided pairs lends credibility to our solution.
\end{itemize}

\clearpage

\section{Problem 2}

\clearpage

\section{Problem 3}

(10 pts) Write up your own conditional example, like “If it is raining, the grass is wet.” but
more interesting and memorable. Write the converse, contrapositive, and the inverse of
your example.

\begin{itemize}
    \item Original Statement: If a student studies hard, then they will pass the exam.
    \item Converse: If a student passes the exam, then they studied hard.
    \item Contrapositive: If a student does not pass the exam, then they did not study hard.
    \item Inverse: If a student does not study hard, then they will not pass the exam.
\end{itemize}

\clearpage

\section{Problem 4}

(10pts) Construct the Truth Tables from Rosen page 15 \#31.

\clearpage

\section{Problem 5}

5. (15pts) Imagine you are given 3 matchboxes, labeled “Red and White,” “Red and Red,”
and “White and White.” Each box contains 2 marbles which can be either red or white.
The labels on each box are incorrect (although the labels would be correct if you
switched them around). You are allowed to peek inside one box and look at exactly one
marble in that box \- and see if it is red or white.

1. After this, can you figure out what is in each box and fix the labels to be correct?

2. Does it matter which box you choose to open at the start?

3. Are you using inductive or deductive reasoning?

4. Can you be sure of your answer?

\clearpage

\section{Problem 6}

\clearpage

\section{Problem 7}

(10 pts) Let’s do a simple deductive proof that a negative number times a negative
number is a positive number.

\clearpage

\section{Problem 8}

(10 pts) Compare the proof in the previous question to the car example video in Moodle.
Both show “a negative times a negative is a positive.”

\clearpage

\end{document}
